% Report di avanzamento: fasi 1-4 della pipeline di compressione sicura basata su BWT.
% Si compila con pdflatex (o latexmk -pdf). Su Overleaf basta incollare l'intero file
% in un progetto nuovo; ogni \section è autonoma e si può copiare anche separatamente.

\documentclass[11pt,a4paper]{article}

\usepackage[T1]{fontenc}
\usepackage[utf8]{inputenc}
\usepackage[italian]{babel}
\usepackage{lmodern}
\usepackage{microtype}
\usepackage[margin=2.5cm]{geometry}
\usepackage{booktabs}
\usepackage{array}
\usepackage{tikz}
\usetikzlibrary{arrows.meta,positioning}
\usepackage[hidelinks]{hyperref}

\setlength{\parskip}{0.4em}

\title{Ottimizzazione e rafforzamento di una pipeline\\di compressione sicura basata su BWT\\[0.5em]
\large Report di avanzamento: fasi prestazionali}
\author{Filippo NF22500049 \and Mattia Gallucci NF22500127 \and Sebastiano NF22500125}
\date{26 novembre 2026}

\begin{document}

\maketitle

\begin{abstract}
Questo report descrive il lavoro svolto sulla pipeline di compressione sicura basata su
sBWT e bMTF. Partendo da un'implementazione interamente in Python, lenta e fragile sui file binari, abbiamo sostituito il calcolo del
suffix array con una libreria nativa, riscritto gli stadi di codifica sui simboli interi,
portato l'intera pipeline a lavorare su byte con un formato binario validato, e reso più
veloce la decompressione. Sui file di testo di riferimento la compressione è oggi circa
quattro volte più veloce, i file compressi sono più piccoli del 39\% e la pipeline accetta
qualunque tipo di file. La fase dedicata al rafforzamento crittografico, è ancora da
svolgere: le debolezze di sicurezza dello schema originale sono quindi ancora presenti e
vengono elencate alla fine del documento.
\end{abstract}

% ---------------------------------------------------------------------------
\section{Introduzione}

Il progetto parte da una pipeline di compressione sviluppata in lavori precedenti, che
combina la compressione con una forma di protezione dei dati.
La versione di partenza funzionava, ma aveva limiti evidenti. Era molto lenta, perché gli
algoritmi più costosi erano scritti in Python puro; leggeva i file come testo e quindi
falliva su qualunque file binario; scriveva su disco diversi file intermedi con i dati in
chiaro; salvava il risultato con un meccanismo di serializzazione che permette di eseguire
codice arbitrario aprendo un file manipolato.

Il lavoro è stato organizzato in un piano a fasi, con alcuni principi condivisi:
\begin{itemize}
  \item \textbf{restare in Python}, spostando le parti lente su librerie native già
        esistenti o su codice compilato con Numba, senza scrivere moduli C propri;
  \item \textbf{misurare prima di ottimizzare}: ogni miglioramento è dimostrato da un
        confronto prima/dopo fatto sempre con lo stesso protocollo;
  \item \textbf{verificare a ogni passo}: finché il formato non cambia l'output deve essere
        identico a quello originale, e quando cambia si verifica che la decompressione
        restituisca esattamente il file di partenza;
  \item \textbf{non inventare crittografia}: lo schema sBWT + bMTF resta quello di
        partenza, cambia il modo in cui è implementato e protetto.
\end{itemize}

\section{Punto di partenza e metodo di misura}

Prima di modificare il codice abbiamo misurato la versione originale (fase 0), per avere
un riferimento con cui confrontare ogni passo successivo. Le misure usano i corpus
standard per la compressione senza perdita (Canterbury, Calgary e, per i file più grandi,
Large e Silesia) e, come riferimento esterno, \texttt{bzip2 -9}, che usa la stessa famiglia
di trasformazioni senza la parte di sicurezza.

Per ogni file si misurano il tempo di compressione e di decompressione, la memoria
massima e il rapporto di compressione, e si verifica che il file decompresso sia identico
all'originale. Ogni prova gira in un processo separato; salvo dove indicato si eseguono tre
ripetizioni più una di riscaldamento, e si riporta la mediana. Distinguiamo due tempi: il
\emph{tempo totale}, che include l'avvio dell'interprete Python e il caricamento delle
librerie, e il \emph{lavoro interno}, cioè il tempo speso nella pipeline vera e propria.

Il confronto principale riguarda i 17 file di testo di Canterbury e Calgary, gli unici che
la versione originale riusciva a elaborare: gli altri 8 file dei due corpus sono binari e
la pipeline originale si interrompeva leggendoli.

% ---------------------------------------------------------------------------
\section{Fase 1 --- Suffix array con una libreria nativa}

Il calcolo della BWT richiede di ordinare tutti i suffissi del testo; il risultato di
questo ordinamento si chiama \emph{suffix array}. Nella pipeline originale questo passo
era svolto da DC3, un algoritmo lineare ma scritto in Python puro, e da solo occupava la
quasi totalità del tempo di compressione.

La soluzione scelta è stata quella di usare una libreria C già esistente, divsufsort.
L'ostacolo apparente era che la sBWT non ordina i simboli nell'ordine naturale ma in un
ordine segreto deciso dalla chiave, mentre una libreria standard conosce solo l'ordine
numerico dei byte. Passandole il testo così com'è si ottiene una BWT normale, invertibile
da chiunque, e la protezione scompare. È probabile che sia stato questo il motivo per cui
nei lavori precedenti la libreria era stata scartata.

L'osservazione chiave è che non serve cambiare l'ordine usato dalla libreria: basta
cambiare il testo. Prima del calcolo ogni simbolo viene sostituito dalla sua posizione
nell'ordine segreto. A quel punto ordinare i numeri in modo crescente equivale a ordinare i
simboli nell'ordine segreto, e la libreria calcola il suffix array segreto senza sapere che
esiste una chiave. Questa rinomina era già presente nel codice, come primo passo di DC3:
è bastato accorgersi che dopo di essa si poteva usare qualunque algoritmo.

\paragraph{Verifica.} Il criterio scelto è più forte del semplice round-trip: a parità di
chiave, la trasformata prodotta con la libreria deve essere identica, simbolo per simbolo,
a quella prodotta con DC3. Così l'output non cambia e quindi non cambia neanche la
sicurezza dello schema. La condizione è stata verificata su tutti i file dei corpus.

\paragraph{Risultati.} La sola sBWT è diventata da 24 a 79 volte più veloce, con un
guadagno che cresce con la dimensione del file; considerando solo il calcolo del suffix
array si arriva a 140--160 volte. Anche la memoria scende molto (su un file di 4~MB da circa
1~GB a 280~MB), perché DC3 rappresentava ogni intero come un oggetto Python. Sull'intera
pipeline la compressione è risultata circa 3,4 volte più veloce: il guadagno è minore
perché, una volta eliminato il suffix array come collo di bottiglia, il tempo si sposta
sugli altri stadi. DC3 è rimasto nel codice come termine di confronto.

\begin{table}[h]
\centering
\begin{tabular}{lrrrr}
\toprule
File & Dimensione & sBWT con DC3 & sBWT con divsufsort & Rapporto \\
\midrule
alice29.txt & 152 KB & 1,47 s & 0,060 s & 24$\times$ \\
book1       & 769 KB & 5,02 s & 0,116 s & 43$\times$ \\
world192.txt& 2,5 MB & 20,6 s & 0,295 s & 70$\times$ \\
bible.txt   & 4,0 MB & 36,5 s & 0,464 s & 79$\times$ \\
\bottomrule
\end{tabular}
\caption{Tempo della sola sBWT su file intero, prima e dopo la fase 1.}
\end{table}

% ---------------------------------------------------------------------------
\section{Fase 2 --- Codifica sui simboli interi}

Dopo la fase 1 il problema principale non era più la velocità ma la dimensione dei file.
Nella versione originale la sequenza prodotta dalla MTF veniva convertita in una stringa di
testo, con cifre e separatori, e la codifica di Huffman lavorava sui singoli caratteri di
quella stringa. In questo modo la struttura statistica creata da BWT e MTF andava in gran
parte persa, e i caratteri di formattazione aggiungevano dati inutili.

La fase 2 ha riscritto i tre stadi successivi alla BWT seguendo l'approccio di bzip2:
\begin{itemize}
  \item la \textbf{MTF} lavora su array di dimensione fissa invece che su liste Python;
        il ciclo, sequenziale per natura, può essere compilato con Numba;
  \item la \textbf{RLE} codifica solo le sequenze di zeri, che dopo BWT e MTF sono la
        grande maggioranza delle ripetizioni, con una rappresentazione binaria compatta che
        richiede un numero di simboli logaritmico nella lunghezza della sequenza;
  \item la \textbf{codifica di Huffman} è stata scritta da noi in forma \emph{canonica}:
        lavora direttamente sui simboli interi e per ricostruire i codici basta conoscerne
        le lunghezze, senza salvare l'albero né usare serializzazioni di Python.
\end{itemize}

L'obiettivo del piano era una riduzione della dimensione di circa il 39\%. È stato
raggiunto: misurata sulla pipeline completa (dopo la fase 3), la dimensione compressa dei
17 file di testo scende da 1.743.298 a circa 1.063.000 byte, il 39\% in meno. Su
\texttt{alice29.txt} il rapporto tra dimensione compressa e originale passa da 0,501 a
0,310, contro lo 0,284 di bzip2.

% ---------------------------------------------------------------------------
\section{Fase 3 --- Pipeline a byte e formato binario}

I singoli componenti delle fasi precedenti erano corretti, ma il programma vero e proprio
era ancora quello originale: leggeva il file come testo, aggiungeva un simbolo sentinella
per marcare la fine dei dati, scriveva nove file intermedi in chiaro e salvava il risultato
con \texttt{pickle}. La fase 3 ha costruito la nuova pipeline mettendo insieme i pezzi.

\paragraph{Tutto in memoria e su byte.} I dati vengono letti come byte e non sono più
interpretati come testo, quindi la pipeline accetta qualunque file. Al posto della
sentinella, che sottraeva un simbolo ai dati ammessi, per ogni blocco si salva l'\emph{indice
primario}, come fa bzip2. L'unico file scritto è quello compresso, e la sua assenza di file
intermedi è verificata automaticamente a ogni esecuzione dei benchmark.

\paragraph{Un formato binario validato.} Il risultato è un unico contenitore con un header a
campi di tipo e dimensione noti, che sostituisce del tutto \texttt{pickle}. L'header prevede
già lo spazio per il nonce e per il tag di autenticazione per la fase di sicurezza successiva e un campo che
identifica il codificatore, così da poterne aggiungere altri in futuro senza rendere
illeggibili i file esistenti. Ogni compressione usa un nonce casuale, per cui lo stesso
file compresso due volte con la stessa chiave produce risultati diversi.

\paragraph{Validazione in decompressione.} Un file compresso arriva dall'esterno e può essere
troncato o manipolato. La regola adottata è che un file non valido deve essere rifiutato
con un errore di formato: mai bloccare il programma, mai farlo leggere fuori dalla memoria
prevista, mai fargli allocare memoria senza limite. Ogni stadio controlla il proprio input:
la coerenza tra i campi dell'header, la correttezza della tabella di Huffman (che deve
descrivere un codice completo e senza sovrapposizioni), la lunghezza delle sequenze di zeri
prima di allocarle, i limiti degli indici nella MTF. Per misurare l'efficacia di questi
controlli abbiamo alterato uno alla volta tutti i byte di un file compresso: l'81\% delle
alterazioni viene rifiutato e nessuna provoca un crash o un blocco. Le restanti producono
dati validi ma diversi: senza un meccanismo di autenticazione non si possono riconoscere per
costruzione, ed è proprio il compito della fase di sicurezza che faremo in seguito.

\paragraph{Problemi trovati lungo il percorso.} Mettendo insieme i componenti sono emersi
alcuni problemi che i test sui singoli moduli non mostravano:
\begin{itemize}
  \item la scrittura dei bit di Huffman aveva un costo quadratico nel numero di simboli, a
        causa di un accumulatore che non veniva mai svuotato: su un file di 769~KB la codifica
        impiegava 13,4~s, ridotti a 0,08~s con una correzione minima e senza cambiare l'output;
  \item i processi paralleli usati per i blocchi potevano bloccarsi definitivamente, per
        un'interazione nota tra la creazione dei processi per copia e le librerie native già
        avviate; ora i processi vengono avviati da zero;
  \item con il suffix array nativo il parallelismo costava più di quanto rendeva sui blocchi
        piccoli: i blocchi sono ora di dimensione fissa (900~KB, come bzip2 -9) e il
        parallelismo si attiva solo sui file sopra i 4~MB;
  \item caricare Numba a ogni avvio rendeva ogni esecuzione circa 3,3 volte più lenta, perché
        il codice veniva ricompilato ogni volta; ora Numba si attiva solo sui dati molto
        grandi e conserva il codice compilato su disco.
\end{itemize}

\paragraph{Scelta del codificatore.} La pipeline originale offriva anche un codificatore
aritmetico e uno LZW. Dalle misure della fase 0, LZW comprime peggio di Huffman su tutti i
file, perché dopo BWT e MTF le ripetizioni che sfrutterebbe sono già state eliminate;
l'aritmetico guadagna circa il 2\% ma decomprime 8--10 volte più lentamente. La nuova
pipeline usa quindi solo Huffman; il campo nel formato lascia aperta la possibilità di
modelli più ricchi in futuro.

\paragraph{Risultati.} Tutti i 25 file di Canterbury e Calgary vengono ora compressi e
decompressi correttamente, contro i 17 della versione originale. La compressione resta ai
livelli della fase 1, mentre la decompressione era diventata più lenta (da 3,29 a 3,87~s
sui 17 file di testo), per il nuovo formato, i controlli aggiunti e la decodifica ancora in
Python. Recuperare questo ritardo era l'obiettivo della fase 4.

% ---------------------------------------------------------------------------
\section{Fase 4 --- Decompressione veloce}

La decompressione aveva tre stadi che procedevano un simbolo alla volta in Python: la
decodifica di Huffman, la MTF inversa e la BWT inversa. La fase è stata divisa in tre
interventi, uno per stadio, misurando dopo ciascuno.

\paragraph{Decodifica di Huffman.} Invece di leggere un bit alla volta e scendere
nell'albero dei codici, il decodificatore legge un gruppo di bit e trova il simbolo con una
tabella precalcolata; solo i codici più lunghi passano ancora dall'albero. Durante il lavoro
è emerso nel lettore di bit lo stesso difetto quadratico già corretto nello scrittore.

\paragraph{BWT inversa.} Invertire la BWT significa seguire una catena di posizioni, in cui
ogni passo dipende dal precedente: per un blocco da 900~KB sono 900.000 passi, che in Python
pesano molto. In un primo passo il calcolo della mappa che lega i passi della catena è stato
reso vettoriale con NumPy. In un secondo passo abbiamo applicato la tecnica del \emph{pointer
jumping}: si costruiscono salti di 2, 4, 8, \dots{} passi, ciascuno ottenuto con una sola
operazione vettoriale, e con circa venti operazioni si ricostruisce l'intero blocco. Il
lavoro totale cresce leggermente, ma si svolge tutto dentro la libreria nativa. Lo stadio è
passato da 0,98 a 0,23~s sul benchmark completo, e la memoria per blocco è scesa da 52 a
14~MB. Le posizioni visitate sono esattamente le stesse di prima, quindi l'output è identico
anche su dati corrotti, e il numero di operazioni dipende solo dalla lunghezza del blocco,
non dal suo contenuto.

\paragraph{MTF inversa.} L'ultima modifica produce direttamente i byte di uscita, invece di
passare da liste di interi Python, ed evita una seconda scansione completa dell'output per
controllarlo, sostituita da un controllo equivalente sul dizionario. Lo stadio è sceso del
38\%. Il controllo dei limiti sui valori in ingresso resta prima del ciclo: abbiamo
verificato che, togliendolo, la versione compilata legge fuori dalla memoria e il processo
va in crash.

Per questa fase il piano prevedeva Numba. Abbiamo invece usato NumPy e tabelle di lookup,
perché sui file del benchmark il costo di avvio di Numba supera il guadagno; Numba resta in
uso solo sui dati molto grandi.

\begin{table}[h]
\centering
\begin{tabular}{lrr}
\toprule
Momento & Tempo totale & Lavoro interno \\
\midrule
Fase 1 (riferimento)          & 3,29 s & --- \\
Inizio fase 4 (dopo la fase 3)& 3,87 s & 2,61 s \\
Dopo la decodifica di Huffman & 3,30 s & 2,14 s \\
Dopo la BWT inversa           & 3,19 s & 1,69 s \\
Dopo la MTF inversa           & \textbf{3,09 s} & \textbf{1,54 s} \\
\bottomrule
\end{tabular}
\caption{Decompressione dei 17 file di testo durante la fase 4.}
\end{table}

Il lavoro interno è sceso del 41\% e la decompressione è tornata sotto il livello della
fase 1. Il tempo totale scende meno, perché ormai circa metà di esso è l'avvio di ogni
processo (interprete e librerie, circa 0,09~s per file).

% ---------------------------------------------------------------------------
\section{Risultati complessivi}

La tabella~\ref{tab:complessivo} confronta la versione originale con quella attuale sugli
stessi 17 file di testo (3,46~MB in tutto).

\begin{table}[h]
\centering
\begin{tabular}{lrrl}
\toprule
Misura & Fase 0 & Oggi & Variazione \\
\midrule
Compressione (tempo totale)      & 15,91 s & 3,89 s & 4,1$\times$ più veloce \\
Compressione (lavoro interno)    & 14,21 s & 2,43 s & 5,8$\times$ più veloce \\
Decompressione (tempo totale)    & 3,22 s  & 3,09 s & $-4\%$ \\
Decompressione (lavoro interno)  & 2,44 s  & 1,54 s & $-37\%$ \\
Dimensione compressa             & 1.743.298 B & 1.062.942 B & $-39\%$ \\
File supportati (Canterbury e Calgary) & 17 su 25 & 25 su 25 & binari compresi \\
\bottomrule
\end{tabular}
\caption{Versione originale e versione attuale a confronto.}
\label{tab:complessivo}
\end{table}

Rispetto a bzip2 -9 i file compressi sono più grandi dell'8,8\% (nella versione originale
erano più grandi del 78,5\%), mentre i tempi restano circa dieci volte superiori. La
differenza di velocità si spiega con il fatto che bzip2 è scritto interamente in C, mentre
noi restiamo in Python con librerie native, e con il lavoro che bzip2 non fa: l'ordinamento
segreto, la permutazione per blocchi della MTF e i controlli contro i file manipolati.

Come verifica finale prima di questo report abbiamo eseguito l'intera suite di test (97
test, tutti superati) e un ciclo completo di compressione e decompressione su tutti e
quattro i corpus, compresi i file binari e i file Silesia fino a 10~MB: tutte le 36 prove
hanno restituito esattamente il file originale.

\section{Prossimo passo: Rafforzamento}

La fase 5, non ancora iniziata, affronta le debolezze elencate sopra usando solo primitive
standard e librerie consolidate:
\begin{itemize}
  \item una \textbf{gerarchia di chiavi}: dalla password e dal nonce del file, tramite una
        funzione di derivazione lenta pensata per le password, si ottiene una chiave
        principale, da cui si ricavano chiavi separate per le permutazioni e per
        l'autenticazione, e infine una chiave per ogni blocco, distinta tra sBWT e bMTF;
  \item \textbf{permutazioni} generate da un flusso pseudocasuale crittografico, al posto
        sia del generatore attuale sia della catena di MD5, così che ogni blocco sia
        indipendente;
  \item uno schema \textbf{Encrypt-then-MAC}: un tag di autenticazione su header e dati,
        verificato per primo in decompressione, prima di decodificare qualunque cosa.
\end{itemize}
\end{document}
